\exercisegroup

\begin{enumerate}
\def\labelenumi{\arabic{enumi})}
\tightlist
\item
  Let \((\mathrm{C}_{i})_{i\in\mathrm{I}}\) be a family of categories. Let C be the sum quiver of the family of underlying quivers of the categories \(C_{i}\).
\end{enumerate}

\begin{enumerate}
\def\labelenumi{\alph{enumi})}
\item
  Show that the quiver C, equipped with the composition law induced by the composition laws of the categories \(C_{i}\), is a category. We call it the sum category of the family \(C_{i}\).
\item
  Prove that the canonical morphism from \(C_{i}\) to C is a functor.
\item
  If the \(C_{i}\) are groupoids, prove that the category C is a groupoid.
\end{enumerate}

\begin{enumerate}
\def\labelenumi{\arabic{enumi})}
\setcounter{enumi}{1}
\tightlist
\item
  Let \((\mathrm{C}_{i})_{i\in\mathrm{I}}\) be a family of categories. Let C be the product of the underlying quivers of the categories \(C_{i}\).
\end{enumerate}

\begin{enumerate}
\def\labelenumi{\alph{enumi})}
\item
  Show that the quiver C, equipped with the composition law induced by the composition laws of the categories \(C_{i}\), is a category. It is called the product category of the family \((\mathrm{C}_{i})\).
\item
  Prove that the canonical morphism \(pr_{i}: C \rightarrow C_{i}\) is a functor.
\item
  If the \(C_{i}\) are groupoids, prove that the category C is a groupoid.
\end{enumerate}

\begin{enumerate}
\def\labelenumi{\arabic{enumi})}
\setcounter{enumi}{2}
\tightlist
\item
  Let C be a category; denote its composition law by m. Let C° be the opposite directed graph equipped with the composition law given by \(m^{\circ}(f,g)=m(g,f)\) for every pair \((f,g)\) of composable arrows of C.
\end{enumerate}

\begin{enumerate}
\def\labelenumi{\alph{enumi})}
\item
  Prove that C° is a category. It is called the opposite category of C.
\item
  If C is a groupoid, prove that C° is a groupoid isomorphic to C.
\end{enumerate}

\begin{enumerate}
\def\labelenumi{\arabic{enumi})}
\setcounter{enumi}{3}
\tightlist
\item
  We define the mass \(\mu(\mathrm{G})\) of a groupoid G as the supremum of the sums \(\sum_{a\in\mathrm{A}}\mathrm{Card}(\mathrm{G}_{a})^{-1}\), where A ranges over the set of subsets of \(\mathrm{Som}(\mathrm{G})\) which meet each orbit of G in at most one point. We say that G has finite mass if \(\mu(\mathrm{G})\in\mathbf{R}\); we say that G is essentially finite if the set of its orbits is finite and if the isotropy group at each of its points is finite.
\end{enumerate}

\begin{enumerate}
\def\labelenumi{\alph{enumi})}
\item
  Show that an essentially finite groupoid has finite mass.
\item
  Let G be the groupoid associated with a finite group \(\Gamma\). It is essentially finite, and we have \(\mu(\mathrm{G}) = 1/\mathrm{Card}(\Gamma)\).
\item
  Let \(\Gamma\) be a finite group, let X be a finite set equipped with a right action of \(\Gamma\), and let G be the associated groupoid. It is essentially finite, and we have \(\mu(\mathrm{G}) = \mathrm{Card}(\mathrm{X}) / \mathrm{Card}(\Gamma)\).
\end{enumerate}

\(d)\) Let \((\mathrm{G}_i)_{i\in \mathrm{I}}\) be a family of groupoids of finite mass. The sum groupoid and the product groupoid of the family \((\mathrm{G}_i)\) have mass \(\sum \mu (\mathrm{G}_i)\) and \(\prod \mu (\mathrm{G}_i)\), respectively.

\begin{enumerate}
\def\labelenumi{\alph{enumi})}
\setcounter{enumi}{4}
\tightlist
\item
  Let X be a set. We define a groupoid G whose set of vertices is the set of finite subsets of X and such that, for every pair \((Y_{1}, Y_{2})\) of finite subsets of X, the set of arrows with origin \(Y_{1}\) and term \(Y_{2}\) is the set of bijections from \(Y_{1}\) to \(Y_{2}\), with composition of arrows given by composition of maps. Show that G is essentially finite and compute its mass.
\end{enumerate}

\begin{enumerate}
\def\labelenumi{\arabic{enumi})}
\setcounter{enumi}{4}
\tightlist
\item
  \begin{enumerate}
  \def\labelenumii{\alph{enumii})}
  \tightlist
  \item
    Let \(\varphi\colon G\to G'\) be a groupoid morphism such that the map \(\operatorname{Som}(\varphi)\) is injective. Prove that \(\varphi(G)\) is a subgroupoid of \(G'\).
  \end{enumerate}
\end{enumerate}

\begin{enumerate}
\def\labelenumi{\alph{enumi})}
\setcounter{enumi}{1}
\tightlist
\item
  Let G be a nontransitive groupoid whose vertex set \(\operatorname{Som}(G)\) is the two-element set \(\{a,b\}\). Let G' be the groupoid associated with the free product \(G_a*G_b\). There exists a unique groupoid morphism \(\varphi:G\to G'\) such that the maps \(\varphi_a\) and \(\varphi_b\) are the canonical maps from \(G_a\) and \(G_b\) to \(G_a*G_b\). If \(G_a\) or \(G_b\) is not the trivial group, the image \(\varphi(G)\) of this groupoid morphism is not a subgroupoid of \(G'\).
\end{enumerate}

\begin{enumerate}
\def\labelenumi{\arabic{enumi})}
\setcounter{enumi}{5}
\tightlist
\item
  Let G and G' be graphs, and let \(\varphi\colon\mathrm{G}'\to\mathrm{G}\) be a graph morphism. We say that \((\mathrm{G}',\varphi)\) is a covering of G if, for every vertex \(a\in\mathrm{Som}(\mathrm{G}')\) and every arrow \(f\in\mathrm{Fl}(\mathrm{G})\) with origin \(\mathrm{Som}(\varphi)(a)\), there exists a unique arrow \(f'\in\mathrm{Fl}(\mathrm{G}')\) with origin \(a\) such that \(\mathrm{Fl}(\varphi)(f')=f\).
\end{enumerate}

We assume that \((\mathrm{G}^{\prime},\varphi)\) is a covering of G.

\begin{enumerate}
\def\labelenumi{\alph{enumi})}
\item
  Prove that there exists a unique action of \(\operatorname{Grp}(G)\) on \(\operatorname{Som}(G')\), relative to the map \(\operatorname{Som}(\varphi)\), such that \(x \cdot \operatorname{Fl}(\varphi)(f) = t(f)\) for every vertex \(x \in \operatorname{Som}(G')\) and every arrow \(f \in \operatorname{Fl}(G')\) with origin x.
\item
  Deduce that if x and y are vertices of G belonging to the same connected component of G, then \(\text{Card}(\text{Som}(\varphi)^{-1}(x)) = \text{Card}(\text{Som}(\varphi)^{-1}(y))\).
\end{enumerate}

\begin{enumerate}
\def\labelenumi{\arabic{enumi})}
\setcounter{enumi}{6}
\tightlist
\item
  Let G be a graph and let \((G_{1},\varphi_{1})\), \((G_{2},\varphi_{2})\) be coverings of the graph G.
\end{enumerate}

\begin{enumerate}
\def\labelenumi{\alph{enumi})}
\item
  Let \(\psi\colon\mathrm{G}_{1}\to\mathrm{G}_{2}\) be a graph morphism such that \(\varphi_{2}\circ\psi=\varphi_{1}\). Show that \(\psi(x\cdot\gamma)=\psi(x)\cdot\gamma\) for all \(x\in\mathrm{Som}(\mathrm{G}_{1})\) and every arrow \(\gamma\) of \(\varpi_{\mathrm{G}}\) with origin \(\varphi_{1}(x)\).
\item
  Let \(u \colon \operatorname{Som}(\mathrm{G}_1) \to \operatorname{Som}(\mathrm{G}_2)\) be a map such that \(u(x \cdot \gamma) = u(x) \cdot \gamma\) for all \(x \in \operatorname{Som}(\mathrm{G}_1)\) and every arrow \(\gamma\) of \(\varpi_{\mathrm{G}}\) with origin \(\varphi_1(x)\). Prove that there exists a unique morphism of graphs \(\psi \colon \mathrm{G}_1 \to \mathrm{G}_2\) such that \(\varphi_2 \circ \psi = \varphi_1\) and \(\operatorname{Som}(\psi) = u\).
\item
  Let X be a set and let \(p: X \to \text{Som}(G)\) be a map. Suppose that an operation of \(\varpi_{G}\) on X is given with respect to p. Prove that there exists a graph \(G'\) with vertex set X and a morphism of graphs \(\varphi: G' \to G\) such that \(\text{Som}(\varphi) = p\), such that \((G', \varphi)\) is a covering of G and the operation of \(\varpi_{G}\) on \(\text{Som}(G')\) coincides with the given operation.
\end{enumerate}

\begin{enumerate}
\def\labelenumi{\arabic{enumi})}
\setcounter{enumi}{7}
\tightlist
\item
  Let G be a connected graph and let \(a\) be a vertex of G.
\end{enumerate}

\begin{enumerate}
\def\labelenumi{\alph{enumi})}
\tightlist
\item
  Let \((\mathrm{G}_{1},\varphi_{1})\) and \((\mathrm{G}_{2},\varphi_{2})\) be coverings of the graph G. Let \(u\colon\operatorname{Som}(\mathrm{G}_{1})\to\operatorname{Som}(\mathrm{G}_{2})\) be a map such that \(u(x)\cdot\gamma=u(x\cdot\gamma)\)
\end{enumerate}

for all \(\gamma \in \varpi_{\mathrm{G},a}\) and every vertex \(x\) of \(\mathrm{G}_1\) such that \(\operatorname{Som}(\varphi_1)(x) = a\). Prove that there exists a unique morphism of graphs \(\psi: \mathrm{G}_1 \to \mathrm{G}_2\) such that \(\psi = \operatorname{Som}(u)\).

\begin{enumerate}
\def\labelenumi{\alph{enumi})}
\setcounter{enumi}{1}
\tightlist
\item
  Let X be a set equipped with an operation of the group \(\varpi_{\mathrm{G},a}\). Prove that there exists a covering \((\mathrm{G}',\varphi)\) of the graph G such that \(\operatorname{Som}(\varphi)^{-1}(a) = \mathrm{X}\) and such that the operation of \(\varpi_{\mathrm{G}}\) on \(\operatorname{Som}(\mathrm{G}')\) induces the given operation of \(\varpi_{\mathrm{G},a}\) on X.
\end{enumerate}

\begin{enumerate}
\def\labelenumi{\arabic{enumi})}
\setcounter{enumi}{8}
\tightlist
\item
  Let G and G' be graphs, and let \(\varphi\colon G'\to G\) be a graph morphism such that \((G',\varphi)\) is a covering of G.
\end{enumerate}

\begin{enumerate}
\def\labelenumi{\alph{enumi})}
\item
  Prove that the graph \(\varphi(\mathrm{Grp}(\mathrm{G}^{\prime}))\) is a subgroupoid of the groupoid \(\mathrm{Grp}(\mathrm{G})\) .
\item
  Prove that, for every vertex \(a\) of \(\mathbf{G}'\), the group homomorphism \(\pi_1(\varphi, a)\) is injective.
\item
  Suppose that G is connected and fix a vertex a of G. Let K be a subgroup of \(\varpi_{G,a}\). Prove that there exists a covering \((H,\psi)\) of the graph G and a vertex b of H such that \(\psi(b)=a\) and K is the image of the morphism \(\pi_{1}(\psi,a)\).
\item
  Prove that a subgroup of a free group is free (see also IV, p. 417, example 2).
\item
  Let n and q be integers; if F is a free group on n generators and K a subgroup of F of index q, prove that K is a free group on \(1 + q(n - 1)\) generators.
\end{enumerate}

\begin{enumerate}
\def\labelenumi{\arabic{enumi})}
\setcounter{enumi}{9}
\tightlist
\item
  Let F be a set, and let C be a subset of \(F \times F\) and \(m: C \to F\) a map. We make the following assumptions:
\end{enumerate}

\begin{enumerate}
\def\labelenumi{(\roman{enumi})}
\item
  The maps from C into F × F given by \((f,g)\mapsto(f,m(f,g))\) and \((f,g)\mapsto(m(f,g),g)\) are injective;
\item
  If \(f, g, h \in \mathrm{F}\) are such that the pairs \((f, g)\) and \((g, h)\) belong to C, then the pairs \((f, m(g, h))\) and \((m(f, g), h)\) belong to C and one has \(m(f, m(g, h)) = m(m(f, g), h)\).
\item
  For every \(f \in \mathrm{F}\), there exist elements \(o(f)\), \(t(f)\), and \(\overline{f} \in \mathrm{F}\), necessarily unique by virtue of (i), such that the pairs \((f, t(f))\), \((o(f), f)\), \((f, \overline{f})\), and \((\overline{f}, f)\) belong to C and one has \(m(f, t(f)) = f = m(o(f), f)\) and \(m(f, \overline{f}) = o(f)\).
\end{enumerate}

\begin{enumerate}
\def\labelenumi{\alph{enumi})}
\item
  Prove that, for every \(f \in F\) , we have \((\overline{f}, f) \in C\) and \(m(\overline{f}, f) = t(f)\) . Deduce that \((o(f), o(f)) \in C\) and \(m(o(f), o(f)) = o(f)\) . Also prove that \(o(\overline{f}) = t(f)\) and \(\overline{\overline{f}} = f\) .
\item
  Let S be the set of elements \(e\) of F such that \((e,e)\in\mathrm{C}\) and \(m(e,e)=e\); let us denote by \(o\) and \(t\) the maps from F to S given by \(f\mapsto o(f)\) and \(f\mapsto t(f)\).
\end{enumerate}

Prove that the quadruple \(\Gamma = (S, F, o, t)\) is a quiver and the map \(m\) is a composition law in \(\Gamma\) which makes it a groupoid. (This construction is due to H. BRANDT; cf. “Über eine Verallgemeinerung des Gruppenbegriffes”, Mathematische Annalen (1926), vol. 96, p. 360–366.)

\begin{enumerate}
\def\labelenumi{\arabic{enumi})}
\setcounter{enumi}{10}
\tightlist
\item
  Let G be a groupoid; let C be the set of composable pairs of arrows of G, and let \(m\colon\mathrm{C}\to\mathrm{Fl}(\mathrm{G})\) be the composition law of G.
\end{enumerate}

\begin{enumerate}
\def\labelenumi{\alph{enumi})}
\item
  Show that C and m satisfy hypotheses (i), (ii), (iii) of Exercise 10. Denote by Γ the groupoid provided by the construction in that exercise.
\item
  Prove that the pair consisting of the map \(e \mapsto o(e)\) from \(\operatorname{Som}(\Gamma)\) to \(\operatorname{Som}(\mathrm{G})\) and of the identity mapping of the set \(\operatorname{Fl}(\mathrm{G})\) into itself is an isomorphism of groupoids from \(\Gamma\) to G.
\end{enumerate}

This shows that a groupoid is determined by the set of its arrows and its composition law.

\begin{enumerate}
\def\labelenumi{\arabic{enumi})}
\setcounter{enumi}{11}
\tightlist
\item
  Let G be a group, let S be the set of subgroups of G, and let F be the set of subsets X of G such that there exist a subgroup H of G and an element g of G with X = gH.
\end{enumerate}

\begin{enumerate}
\def\labelenumi{\alph{enumi})}
\item
  Let \(\mu\) be the map from \(G \times G \times G\) to \(G\) defined by \(\mu(x,y,z)=xy^{-1}z\). Prove that a non-empty subset \(X\) of \(G\) belongs to the set F if and only if \(\mu(X\times X\times X)\subset X\).
\item
  Show that, by setting \(o(\mathrm{X}) = g^{-1}\mathrm{X}\) and \(t(\mathrm{X}) = \mathrm{X}g^{-1}\) for an arbitrary element g of X, one defines maps \(o: F \to S\) and \(t: F \to S\).
\item
  If \(X_{1}\) and \(X_{2}\) are elements of S such that \(t(\mathrm{X}_{1}) = o(\mathrm{X}_{2})\), we set \(m(\mathrm{X}_{1}, \mathrm{X}_{2}) = \mathrm{X}_{1}\mathrm{X}_{2}\). Show that the map m is a composition law in the quiver (S,F,o,t) which makes it a groupoid ("Baer groupoid"). It is denoted \(\mathrm{Ba}(\mathrm{G})\).
\item
  Show that the orbits of \(\mathrm{Ba}(\mathrm{G})\) are the conjugacy classes of subgroups of G.
\item
  Show that the isotropy group of Ba(G) at an element H of S is isomorphic to the quotient group N(H)/H.
\end{enumerate}

\setcounter{exercisegroup}{4}

